\lecture{15}
\title{The Cook-Levin Theorem and Space Complexity}

\begin{document}

\subsection{The Cook-Levin Theorem}

Let's return to the black-boxed result from
\href{/18.404/lectures/14/#np-completeness}{Lecture 14}
and actually provide a proof.

\begin{theorem}{Cook-Levin}
    The language $\mathrm{3SAT}$ is NP-complete.
\end{theorem}

\begin{proof}
    We must show that for any $A \in \mathrm{NP}$, we have
    $A \leq_p \mathrm{3SAT}$.
    Recall from \href{/18.404/lectures/12-13/#polynomial-verifiability-np}{Lecture 12}
    that since $A \in \mathrm{NP}$,
    it follows that there is some NTM $M$
    that decides $A$ in at most $n^k$ steps.

    Any accepting computation history for $M$
    can be represented using at most $n^k$ configurations,
    each with length at most $n^k$.
    Therefore, we can create, in polynomial time, an instance of $\mathrm{3SAT}$ such that:
    \begin{itemize}
        \item There are $n^k \cdot n^k \cdot |\Gamma \cup Q|$ variables
        $x_{i, j, \epsilon}$, where $x_{i, j, \epsilon} = \textsc{True}$
        if and only if the $i^{\text{th}}$ configuration has
        $\epsilon \in \Gamma \cup Q$ in position $j$.
        \item The instance is satisfied
        if and only if the $\{x_{i, j, \epsilon}\}$ specify
        an accepting computation history.
    \end{itemize}
    Details of how this instance looks were provided in lecture,
    but I did not find them very interesting,
    so they are omitted here.

    \includewidget{lecture-15/tableau}

    The (warning: AI-generated) widget above
    shows this reduction for 
    $A = \{w \in \{\str{0}, \str{1}\}^4
    \mid w \text{ contains a } \str{1}\}$,
    where the NTM $M$ simply scans the input string
    and nondeterministically accepts if
    its head is on a $\str{1}$.
    ~ $\blacksquare$
\end{proof}

\subsection{Space Complexity}

\textbf{Definition.}
A TM has \textbf{space complexity} $f(n)$
if it halts on all $w$ of length $n$
using at most $f(n)$ tape cells.
Furthermore:
\begin{itemize}
    \item We denote $\mathrm{SPACE}(f(n)) := \{
        B \mid
        B \text{ is decidable by a TM in space complexity } O(f(n))
    \}$.
    \begin{itemize}
        \item We denote $\mathrm{PSPACE} :=
        \bigcup_k \mathrm{SPACE}(n^k)$.
    \end{itemize}
    \item We denote $\mathrm{NSPACE}(f(n)) := \{
        B \mid
        B \text{ is decidable by an NTM in space complexity } O(f(n))
    \}$.
    \begin{itemize}
        \item We denote $\mathrm{NPSPACE} :=
        \bigcup_k \mathrm{NSPACE}(n^k)$.
    \end{itemize}
\end{itemize}

\begin{theorem}{$\mathrm{TIME} \subseteq \mathrm{SPACE}$}
    For any function $f(n) \geq n$, we have
    $\mathrm{TIME}(f(n)) \subseteq \mathrm{SPACE}(f(n))$.
\end{theorem}

\begin{proof}
    In $f(n)$ time, one can only ever visit
    at most $f(n)$ tape cells. ~ $\blacksquare$
\end{proof}

\begin{theorem}{$\mathrm{SPACE} \subseteq \mathrm{TIME}$}
    For any function $f(n)$, we have
    $\mathrm{SPACE}(f(n)) \subseteq \mathrm{TIME}\left(2^{O(f(n))}\right)$.
\end{theorem}

\begin{proof}
    There are only $2^{O(f(n))}$ configurations,
    so any decider
    cannot possibly take more than $2^{O(f(n))}$ steps,
    or else it would loop and thus never halt.
    ~ $\blacksquare$
\end{proof}

We'll pick up this discussion next lecture.

\end{document}